\section{结论}

张量网络已被提出作为有前景的语言模型，我们首先将 TT 分解应用于真实世界的语言建模数据集，并将该框架命名为 TTLM。我们提出了两个变体：TTLM-Large 和 TTLM-Tiny，并表明它们优于具有低规模隐单元的 Vanilla-RNN。这些实验结果的展示是在机器学习中探索张量网络的一次推进。同时，我们证明 Second-order RNN、RAC 和 MI-RNN 都是 TTLM 的特殊实现。

本研究的一个局限是应当考察不同 normalization 函数的影响。在未来工作中，如果有合适的数学工具和基准测试可用，我们计划研究 TTLM 在自然语言中长期程关联的建模能力，这被认为是 TTLM 的核心特征之一。
% \benyou{no benchmark}
















% In the unusual situation where you want a paper to appear in the
% references without citing it in the main text, use \nocite
% \nocite{langley00}

\bibliography{icml2023}
\bibliographystyle{icml2023}


%%%%%%%%%%%%%%%%%%%%%%%%%%%%%%%%%%%%%%%%%%%%%%%%%%%%%%%%%%%%%%%%%%%%%%%%%%%%%%%
%%%%%%%%%%%%%%%%%%%%%%%%%%%%%%%%%%%%%%%%%%%%%%%%%%%%%%%%%%%%%%%%%%%%%%%%%%%%%%%
% APPENDIX
%%%%%%%%%%%%%%%%%%%%%%%%%%%%%%%%%%%%%%%%%%%%%%%%%%%%%%%%%%%%%%%%%%%%%%%%%%%%%%%
%%%%%%%%%%%%%%%%%%%%%%%%%%%%%%%%%%%%%%%%%%%%%%%%%%%%%%%%%%%%%%%%%%%%%%%%%%%%%%%
